% Work-MAT | 2. rocnik | pripravovana hodina 49
% Návaznost: uživatelské připravy.zip (hodiny 01–20), autorská učebnice na GitHubu.
% Realisticky.cz: didaktická inspirace v posloupnosti a typech úloh; zadání jsou vlastní.
\documentclass[11pt,a4paper]{article}
\input{../../sablony/pripravy/priprava.sty}
\usepackage{tikz}
\nastavPripravu{2. ročník}{Ověřování znalostí}{49}{Průběžný test T4}
\begin{document}
\blok{Motivace a návaznost}
T4 bude krátký test s rozborem; podkladové úlohy jsou analogické, nikoli schválené zadání. Zaměření: exponenciální a logaritmické rovnice, orientovaný úhel a míry.
\poznamka{Orientační nabídka pro 45--55 minut: krátké zopakování a motivace, nový výklad, tři společné vzorové úlohy, samostatná práce. Vyber podle tempa třídy; zbývající zadání slouží jako rezerva.}
\blok{Výklad}
Dbejte na podmínky logaritmů a převody stupně–radiány; v této hodině je rozhodující diagnostika a následná zpětná vazba.
\blok{Společné příklady}
\begin{prikladReseniBox}[title={Příklad 1}]
Vyřeš, resp. urči: $\log_2(x-1)=3$.
\tcblower
\[\begin{aligned}x>1,\quad x-1=8 \\ x=9\end{aligned}\]
\end{prikladReseniBox}
\begin{prikladReseniBox}[title={Příklad 2}]
Vyřeš, resp. urči: $210^\circ\text{ v radiánech}$.
\tcblower
\[\begin{aligned}210^\circ\cdot\frac{\pi}{180^\circ}=\frac{7\pi}{6}\end{aligned}\]
\end{prikladReseniBox}
\poznamka{Ověř aktuální termín, varianty a pravidla klasifikace. Nevytvářej automaticky neodsouhlasené testové zadání.}
\blok{Samostatné procvičování}
\begin{zadaniReseniBox}[title={Samostatná práce 1}]
\begin{enumerate}[label=\alph*)]
\item $765^\circ\text{ po otočkách}$.
\item $-720^\circ\text{ po otočkách}$.
\item $\frac{\pi}{2}\text{ na stupně}$.
\item $\frac{5\pi}{6}\text{ na stupně}$.
\item $\frac{7\pi}{4}\text{ na stupně}$.
\item $45^\circ\text{ na radiány}$.
\end{enumerate}
\tcblower
\begin{enumerate}[label=\alph*)]
\item $45°$.
\item $0°$.
\item $90°$.
\item $150°$.
\item $315°$.
\item $\frac{\pi}{4}$.
\end{enumerate}
\end{zadaniReseniBox}
\blok{Dobrovolné rozšíření}
\begin{zadaniReseniBox}[title={Rozšíření 1}]
\begin{enumerate}[label=\alph*)]
\item Navrhni obdobnou úlohu s jiným úhlem a zkontroluj výsledek graficky..
\end{enumerate}
\tcblower
\begin{enumerate}[label=\alph*)]
\item $Srovnej s jednotkovou kružnicí nebo grafem.$.
\end{enumerate}
\end{zadaniReseniBox}
\blok{Úkol na doma}
Každou hodnotu zdůvodni náčrtkem jednotkové kružnice nebo orientovaným úhlem; u přesných hodnot neuváděj jen desetinné aproximace.
\begin{zadaniReseniBox}[title={Domácí úkol}]
\begin{enumerate}[label=\alph*)]
\item $120^\circ\text{ na radiány}$.
\item $270^\circ\text{ na radiány}$.
\item $330^\circ\text{ na radiány}$.
\item $\frac{11\pi}{6}\text{ na stupně}$.
\item $-90^\circ\text{ na radiány}$.
\item $\frac{3\pi}{4}\text{ na stupně}$.
\item $225^\circ\text{ na radiány}$.
\item $-45^\circ\text{ jako kladný úhel}$.
\end{enumerate}
\tcblower
\begin{enumerate}[label=\alph*)]
\item $\frac{2\pi}{3}$.
\item $\frac{3\pi}{2}$.
\item $\frac{11\pi}{6}$.
\item $330°$.
\item $-\frac{\pi}{2}$.
\item $135°$.
\item $\frac{5\pi}{4}$.
\item $315°$.
\end{enumerate}
\end{zadaniReseniBox}
\poznamka{Domácí úkol obsahuje 8 položek; předpokládaná domácí práce přibližně 15--30 minut. Po větších blocích jsou vědomě promíchané dříve probrané typy.}
\end{document}
